\subsection{Shrinking}
\label{subsection:shrink}
We are about to prove that the sweep rule shrinks finite errors.  
The following claim relates the local maximum points of the potential walls $L_{1}, \ldots, L_{d}$ in an assignment $z$ to the structure of the syndrome they observe. Later, we will use that relation to infer that the decoding step of the sweep rule can't spread the error to a vertex\footnote{To faces whose vertices have higher values of $L_{1}, \ldots, L_{d}$ than the current global maximum values.} with a higher value than the current global maximum of $L_{1}, \ldots, L_{d}$. In other words, the error does not spread in the positive directions.

\begin{claim}
  \label{claim:edge_pole}
Let $z$ and $\xi$ be $d^{\prime}$-dimensional connected finite chains such that $\xi$ results from the application of the sweep rule on $z$.
For any generator $g \in S$, the maximal values of $L_{g}$ in $\xi$ and $\partial^{-} \xi$ are lower than the maximal values in $z$ and $\partial^{-} z$. Namely:
\begin{equation*}
  \begin{split}
    & \max \left\{ L_{g}(u) : u \in \xi \right\}  \le  \max \big\{ L_{g}(u) : u \in z \big\} \\
    & \max \left\{ L_{g}(u) : u \in \partial^{-} \xi \right\}  \le  \max \big\{ L_{g}(u) : u \in \partial^{-} z \big\} 
  \end{split}
\end{equation*}

\end{claim}
\begin{proof}
  First, we argue that for a vertex $v$ in $\partial^{-} z$, and a generator $g \in S$, such that $v$ is a local maximum of $L_{g}$ in either $z$ or $\partial^{-} z$, it implies that $gv \notin \partial^{-} z$. The contradiction where $v$ is a maximum of $L_{g}$ in $z$ follows immediately, since containment in $\partial^{-} z$ implies containment in $z$. Where $v$ is a maximum in $\partial^{-} z$, suppose through contradiction that $gv \in \partial^{-} z$, then
  \begin{equation*}
    \begin{split}
      L_{g}(gv) = (gv)_{g} = 1 + v_{g} = 1 + L_{g}(v)
    \end{split}
  \end{equation*}
In contradiction to $v$ being a local maximum of $L_{g}$ in $\partial^{-} z$.

For a vertex $v \in \partial^{-}z$ denote the correction it applies by $\varphi_{v} \in \mathbb{F}_{2}^{\mathcal{F}_{d^{\prime}}} $, so 
\begin{equation*}
  \begin{split}
    \xi = \sum_{v \in \partial^{-}z}\varphi_{v} + z \ \ \text{ and } \ \ \partial^{-}\xi = \sum_{v \in \partial^{-}z}\partial^{-}\varphi_{v} + \partial^{-}z
  \end{split}
\end{equation*}
So for any vertex $u \in \xi/z$, there is a vertex $v \in \partial^{-}z$ such that $u \in \varphi_{v}$. Moreover, since $\phi_{v}$ is a rooted chain, it follows that $L_{g}(u) \le L_{g}(v) + 1$, where the equality holds when $g$ belongs to the generator's subset of the face $\varphi_{v}$. Similarly for $u \in \partial^{-}\xi/\partial^{-}z$. 

Now, combining with the sweep rule's first condition, which restricts the correction applied by a vertex to have a vertex support that is contained in the positive syndrome it sees, we have that for any vertex $v$ which is a local maximum of $L_{g}$ (in either $z$ or $\partial^{-}z$), any $u \in \varphi_{v}$ satisfies $L_{g}(u) = L_{g}(v)$. In other words, $\xi$ and $\partial^{-}\xi$ don't contain vertices of the form $gv$ where $v$ is a local maximum of $L_{g}$ in $z$ and $\partial^{-}z$, and the desired inequalities follow.
\end{proof}


Next, we prove that for a finite $X$-error, the sweep-cycle advances the corner (the maximum of $L_{d+1}$) of its reduced form. We assume we are given a finite error in its reduced form, and therefore whose corner sees a syndrome, and show that the sweep-rule application excludes the corner. In other words, the sweep-rule shrinks $X$-errors.

\begin{claim}
  \label{claim:sweepshrink}
Let $z$ be a finite reduced $d^{\prime}$-dimensional chain, and let $\xi$ be the reduced result of applying the sweep rule to it. Then:
  \begin{enumerate}
    \item For any generator $g \in S$: $\max \big\{ L_{g}(v) : v \in z \big\} \ge \max \big\{ L_{g}(v) : v \in \xi \big\}  $
    \item And for $L_{d+1}$ : $ \max \big\{ L_{d+1}(v)  : v \in z \big\} \ge \max \big\{  L_{d+1}(v) : v \in \xi \big\} + 1$
  \end{enumerate}
  In particular, $\Size{\xi} \le \Size{z} - 1$.
\end{claim}
\begin{proof}
  The first inequality was proved in \Cref{claim:edge_pole}, so it's left to prove the second. Denote by $\varphi$ the applied correction and by $v$ the vertex on which $L_{d+1}$ gets the maximal value in $z$, if there are more then one vertex then choose $v$ arbitrarily. Since $z$ is reduced, then by \Cref{claim:reducederr} it holds that $v \in \partial^{-} z$. Moreover, since $v$ is a local max of $L_{d+1}$ in $z$ the two follows: 
  \begin{enumerate}
      \item \Cref{claim:findcor} guarantees that $\partial^{-} \varphi|_{v} = \partial^{-} z|_{v}$. Thus, adding $\varphi|_{v}$ to $z$ excludes $v$ from the boundary.
      \item For any other vertex $u \in z$, there is no face $\rotfv{u}{S^{\prime}} \in z$ such that $v \in \rotfv{u}{S^{\prime}}$. Otherwise, it would contradict $v$ being a local maximum of $L_{d+1}$ in $z$. Thus, $v \notin \textbf{Str}_{0+}(u, \partial^{-} z)$, so by \Cref{claim:findcor}, $v \notin \varphi|_{u}$.
  \end{enumerate}
Combining the two, we get that $v \notin \partial^{-} (z + \varphi)$, so the reduced form of $z + \varphi$ does not contain $v$. Namely, $v \notin \xi$. Repeating the above for all the other vertices which get the maximal $L_{d+1}$-value in $z$ yields that $L_{d+1}(\xi) \le L_{d+1}(z) - 1$.

\end{proof}
We finished analyzing the first part of the sweep-cycle that acts in the computational basis. In the next section, we show that the second part of the cycle shrinks reduced phase errors.

\begin{corollary}
\label{claim:polesmove}  
Consider an Pauli $P$, whose $X$- and $Z$-components are set on the reduced chains $x, z \in \mathbb{F}_{2}^{\mathcal{F}_{d^{\prime}}}$. Denote by $\tilde{P}$ the Pauli operator result by an application of the sweep-rule on $P$ and by $\tilde{x}, \tilde{z} \in \mathbb{F}_{2}^{\mathcal{F}_{d^{\prime}}}$ the chains corresponding to the $X$- and $Z$-components of $\tilde{P}$. If both $x, \phi(z)$ are finite connected chain, Then:
  \begin{enumerate}
    \item For any $i \in [d]$: 
      \begin{equation*}
        \begin{split}
          & \max_{v \in x} L_{i}(v)  \ge \max_{v \in \tilde{x}} L_{i}(v) \\
          & \max_{v \in \phi(z)} L_{i}(v)  \ge \max_{v \in \phi(\tilde{z}) } L_{i}(v) 
      \end{split}
      \end{equation*}
    \item And for $i = d+1$ : 
      \begin{equation*}
        \begin{split}
          & \max_{v \in x} L_{d+1}(v)  \ge \max_{v \in \tilde{x}} L_{d+1}(v) + 1 \\
          & \max_{v \in \phi(z)} L_{d+1}(v)  \ge \max_{v \in \phi(\tilde{z}) } L_{d+1}(v) + 1
        \end{split}
      \end{equation*}
  \end{enumerate}
  In particular: 
  \begin{equation*}
    \begin{split}
      &\Size{\tilde{x}}  \le \Size{x}  - 1 \\
      &\Size{ \phi(\tilde{z} ) }  \le \Size{ \phi(z) } - 1 
    \end{split}
  \end{equation*}
\end{corollary}

%\subsubsection{The Sweep-Rule Phase in the Dual Code}
%
%\inb{\textbf{ Eventually we would not need that section for the main theorem. Yet, it might be useful for the Clifford-free case, so do not remove it. }}
%
%
%
%The $Z$-errors are addressed in the second phase of the sweep cycle. In this stage, we move to the dual code, which is a $(d,d-d^{\prime})$-Toric code. So, \Cref{claim:sweepshrink} ensures that the error is reduced relative to the dual image. Yet, what we care about is that the error will be shrunk in the original view, namely after turning back to the $(d,d^{\prime})$-Toric code.
%To do that, we will need the following claim:
%\begin{claim}
%  \label{claim:checksup}
%  Let $\sigma$ be a $d^{\prime}-1$ dimensional chain. Such a vertex $v$ suggests a correction. Denote $w = \sigma|_{v+}$ and for any $g \in S$, denote by $w|_{gv}$ and $w|_{/gv}$ the restriction of $w$ to faces containing $gv$ and their complement. If $w|_{gv} \neq 0$, then $w|_{/gv} \neq 0$.
%\end{claim}
%\begin{proof}
%  It's given that $v$ applies a correction, denoted by $z \in \mathbb{F}_{2}^{\mathcal{F}_{d^{\prime}}}$, and recall that $z$ is rooted in $v$, namely $z = z|_{v} = z|_{v+}$, and $\partial^{-} z|_{v} = w$. By the zeroing of boundary of boundary, we have that $(\partial^{-})^{2} z = 0$. Hence:
%  \begin{equation*}
%    \begin{split}
%      (\partial^{-}w)|_{v+} = (\partial^{-}( (\partial^{-} z) |_{v+}))|_{v+} = (\partial^{-}( \partial^{-} z))|_{v+} = 0  
%    \end{split}
%  \end{equation*}
%Assume that $w|_{/gv} = 0$, which means that there are $d^{\prime}-2$-dimensional faces $A = \{\rotfv{v}{S^{\prime}}\}$ such that $w$ can be written as:
%  \begin{equation*}
%    \begin{split}
%      w = \sum_{  \rotfvs{} \in A  } \rotfv{v}{S^{\prime} \cup g }
%    \end{split}
%  \end{equation*}
%  and therefore:
%  \begin{equation*}
%    \begin{split}
%      (\partial^{-}w)|_{v} =  \sum_{  \rotfvs{} \in A  } \rotfv{v}{S^{\prime}  } + \sum_{  \rotfvs{} \in A  } \sum_{g^{\prime} \in S^{\prime}/g}\rotfv{v}{(S^{\prime} / g^{\prime}) \cup g  } 
%    \end{split}
%  \end{equation*}
%Combining that $(\partial^{-}w)|_{v} = 0$ and that all the right terms have support on $gv$ while the left doesn't, we have that:
%  \begin{equation*}
%    \begin{split}
%      \sum_{  \rotfvs{} \in A  } \rotfv{v}{S^{\prime}  }  = 0
%    \end{split}
%  \end{equation*}
%So any $d^{\prime}-2$-dimensional face $\rotfv{v}{S^{\prime}}$ has to appear an even number of times, but then it follows that any $d^{\prime}-1$-dimensional face of the form $\rotfv{v}{S^{\prime} \cup g }$ has to appear an even number of times in $w$, namely $w = 0$, which is a contradiction to the assumption that $v$ suggests a correction.
%\end{proof}
%
%\begin{claim}
%  \label{claim:corrdual}
%Consider a finite $d^{\prime}$-dimensional chain $z \in \mathbb{F}_{2}^{\mathcal{F}_{d^{\prime}}}$ and assume that its $\phi$-transpose chain $\phi(z)$ is also finite. Denote by $\xi \in \mathbb{F}_{2}^{\mathcal{F}_{d- d^{\prime}}}$ the result of applying the sweep rule over the reduced form of $\phi(z)$. Then, for any generator $g \in S$, we have:
%  \begin{equation*}
%    \begin{split}
%      \max_{v \in \phi^{-1}(\xi)}L_{g}(v) \le \max_{v \in z }L_{g}(z)
%    \end{split}
%  \end{equation*}
%  Furthermore, for the $d+1$th potential wall, we have:
%  \begin{equation*}
%    \begin{split}
%      \max_{v \in \phi^{-1}(\xi)}L_{d+1}(v) \le \max_{v \in z }L_{d+1}(z) - 1
%    \end{split}
%  \end{equation*}
%\end{claim}
%\begin{proof}
%We start by showing the first inequality. Assume, for the sake of contradiction, that this is not the case. In particular, there must be a vertex $u$ in $\phi^{-1}(\xi)$ that is not in $z$. Denote the face containing it by $\rotfvs{}$. Since $u \notin z$, it follows that $\phi(\rotfvs{}) \notin \phi(z)$. Moreover, since $\rotfvs{} \in \phi^{-1}(\xi)$, it follows that $\phi(\rotfvs{}) \in \xi$. Namely, it was 'added' to $\phi(z)$ during the sweep rule.
%
%Let us express explicitly $\phi(\rotfvs{}) = \rotfv{S^{\prime} v}{S/S^{\prime}}$, and since a non-trivial correction is applied, there has to exist a check, rooted in $S^{\prime}v$, that is not satisfied in $\phi(z)$. Denote it by $\rotfv{S^{\prime}v}{J}$ such that $J \subset S/S^{\prime}$. \Cref{claim:checksup} implies that w.l.g. we can assume that $g \notin J$.
%
%
%So $\rotfv{S^{\prime}v}{J} \in \partial^{-}\phi(z)$; therefore, $\rotfv{ (S/J)^{-1} S^{\prime}  v}{S/J} \in \partial^{+}z$. Since $g \notin J$, it follows that $g \in S/J$. Split based on whether $g \in S^{\prime}$. If $g \in S^{\prime}$, then:
%\begin{equation*}
%  \begin{split}
%      L_{g}( (S/J)^{-1} S^{\prime} v) = L_{g}(v)
%  \end{split}
%\end{equation*}
%And hence:
%\begin{equation*}
%  \begin{split}
%    \max_{w \in  \rotfv{  (S/J)^{-1} S^{\prime}  v}{S/J}  }L_{g}(w) = 1 + L_{g}(v) \ge  \max_{w \in  \rotfv{v}{S^{\prime}}  } L_{g}(w) = L_{g}(u)
%  \end{split}
%\end{equation*}
%On the other hand, if $g \notin S$, then:
%\begin{equation*}
%  \begin{split}
%      L_{g}( (S/J)^{-1} S^{\prime} v) = -1 + L_{g}(v)
%  \end{split}
%\end{equation*}
%Yet, for any vertex $w$ in the face $\rotfv{v}{S^{\prime}}$, it holds that $L_{g}(w) = L_{g}(v)$. Thus:
%\begin{equation*}
%  \begin{split}
%    \max_{w \in  \rotfv{   (S/J)^{-1} S^{\prime} v}{S/J}  }L_{g}(w) = 1 - 1 + L_{g}(v) = \max_{w \in  \rotfv{v}{S^{\prime}}  } L_{g}(w) = L_{g}(u)
%  \end{split}
%\end{equation*}
%In other words, we found a vertex in $z$ with an equal or higher $L_{g}$-value than $u$.
%
%For the second part, let us consider the action of $\phi$ on a rooted chain. Let $y$ be a rooted chain at vertex $v$, namely there are faces $A = \{ \rotfvs{} \}$, such that:
%\begin{equation*}
%  \begin{split}
%    y = \sum_{ \rotfvs{} \in A}\rotfvs{}
%  \end{split}
%\end{equation*}
%Thus, 
%\begin{equation*}
%  \begin{split}
%    \phi(y) = \sum_{ \rotfvs{} \in A}\rotfv{S^{\prime}v}{S/S^{\prime}}
%  \end{split}
%\end{equation*}
%So if $v$ achieves the maximal $L_{d+1}$-value in $z$, then the vertices that get the maximal value in $\phi(z)$ are of the form $S^{\prime}v$, for some $d^{\prime}$-size subset of generators $S^{\prime} \subset S$. Therefore, the following equality holds:
%\begin{equation*}
%  \begin{split}
%    \max_{v \in \phi(z)}L_{d+1}(v) = \max_{v \in z }L_{d+1}(v) - d^{\prime}
%  \end{split}
%\end{equation*}
%Now, since the sweep rule decreases the $L_{d+1}$-value (\Cref{claim:sweepshrink}), we have that:
%\begin{equation*}
%  \begin{split}
%    \max_{v \in \xi }L_{d+1}(v) \le \max_{v \in \phi(z) }L_{d+1}(v) - 1
%  \end{split}
%\end{equation*}
%So, in overall:
%\begin{equation*}
%  \begin{split}
%    \max_{v \in \phi^{-1}(\xi) }L_{d+1}(v) & =  \max_{v \in \xi }L_{d+1}(v) + d^{\prime}  \le \max_{v \in \phi(z) }L_{d+1}(v) - 1 + d^{\prime} \\
%    & \le \max_{v \in z }L_{d+1}(v) - 1  
%  \end{split}
%\end{equation*}
%\end{proof}
%Combining \Cref{claim:sweepshrink} and \Cref{claim:corrdual},


%We can state that a sweep-cycle shrinks a finite connected error.
%\begin{corollary}
%\label{claim:polesmoveB}  
%Consider an error whose $X$- and $Z$-product components lie at positions corresponding to the reduced chains $z_{X}, z_{Z} \in \mathbb{F}_{2}^{\mathcal{F}_{d^{\prime}}}$. Denote by $\xi_{X}, \xi_{Z} \in \mathbb{F}_{2}^{\mathcal{F}_{d^{\prime}}}$ chains corresponding to the $X, Z$ positions of the result after the application of a full sweep-cycle. If both $z_{X}, z_{Z}$ are finite connected chain, Then:
%  \begin{enumerate}
%    \item For any $i \in [d]$: 
%      \begin{equation*}
%        \begin{split}
%      \max \big\{ L_{i}(v) : v \in  z_{X} \cup z_{Z} \big\} \ge \max \big\{ L_{i}(v) : v \in  \xi_{X} \cup   \xi_{Z} \big\}  
%        \end{split}
%      \end{equation*}
%    \item And for $i = d+1$ : 
%      \begin{equation*}
%        \begin{split}
%      \max \big\{ L_{d+1}(v)  : v \in  z_{X}  \cup z_{Z} \big\} \ge \max \big\{  L_{d+1}(v) : v \in  \xi_{X} \cup  \xi_{Z}  \big\} + 1
%        \end{split}
%      \end{equation*}
%  \end{enumerate}
%\end{corollary}

\subsubsection{Ideal decoder}
 Although the main role of the section is to show that a single sweep-cycle keeps error clusters under control during a noisy run (by admitting a modest shrinking), an immediate corollary of this is that the application of non-noisy multiple iterations of the sweep-cycle has the potential to erase errors, namely a decoder. Here, a classical measurement is given, we act on classical bits, and the sweep-cycle consists only of the first stage that works over the computational basis, without moving to the dual code. We formulate the corollary as follows:
\begin{corollary}
  \label{col:idealcorrection} 
  Consider the $(d,d^{\prime})$-Toric code, induced by a Toric with side length $n$. Given a connected finite reduced error $z \in \mathbb{F}_{2}^{\mathcal{F}_{d^{\prime}}}$ on that Toric, with $\Size{z} \le n-1$. Then an application of $\Size{z}$ sweep cycles (partial cycles consisting only of the first stage) erases the error.
\end{corollary}
When $z$ is not connected, disjoint components might merge, and it is no longer true that the size of the chain decreases in each iteration. However, we will see later (\Cref{lemma:keylemma}) that performing $\Theta(n^{d+1})$ sweep-cycles erases the error with high probability.



